\section{物理項の改訂と複数ガスへの展開}
\subsection{物理項を変更する三つの方法}
OESCRでは、改訂の内容に応じて、データと設定の変更、関数形を置き換えるpluginの追加、連立方程式を拡張する計算核の開発を使い分ける。日常的なモデル改訂の中心は、状態・反応・遷移・分子帯を記したYAMLと、断面積CSVの更新である。

例えばO 777 nmの消光係数を改訂する作業では、対応する二体過程の係数を更新する。放射分岐の追加では遷移集合を更新し、励起断面積の変更では反応が参照するCSVを差し替える。供給と損失を同じ行列組立て規則で評価するため、これらの改訂は準位寿命と発光へ一貫して反映される。変更前後の準位収支とスペクトルを比較すれば、その物理項が観測へ及ぼす影響を調べられる。

\begin{table}[htbp]
\centering\small
\begin{tabularx}{\linewidth}{p{31mm}p{42mm}Y}
\toprule
改訂内容 & 実装済みの編集箇所 & 改訂時に確認する物理量 \\
\midrule
励起断面積 & 反応のCSVまたは\code{cross\_section\_file} & 状態・閾値・断面積単位・energy範囲・出典 \\
放射寿命・分岐 & species packの\code{transitions} & 全分岐による損失、air/vacuumの波長基準 \\
消光・一次失活 & \code{two\_body}/\code{first\_order}反応 & 衝突相手の入力密度と係数単位 \\
三体過程 & \code{three\_body}反応 & 二つの外部衝突相手の密度、係数の定義 \\
壁損失 & \code{wall}の$\gamma,L,f_{\rm flux}$ & 質量・熱流束定義・有効長 \\
放射捕獲 & 遷移の\code{trapping} & slab脱出係数、入力光学厚さ$\tau$ \\
分子帯 & \code{bands}の率・収率・profile & 光子収率、出力単位、profile支持域 \\
観測条件 & 視線幾何と\code{instruments} & 光路、透過率、LSF、画素、校正基準 \\
関数形 & 対応するplugin registry & 単位、正規化、解析解・保存則 \\
\bottomrule
\end{tabularx}
\caption{物理項ごとの改訂経路。既存の反応形式と観測演算を利用する変更は、該当するデータ・設定へ局所化できる。}
\label{tab:edits}
\end{table}

関数形を変更する場合は、EEDF、reaction rate、band emission/profile、trapping、wall loss、geometry、throughput、LSF、baselineのregistryへpluginを登録する。各pluginは選択キー\code{kind}、説明、入力schema、実行関数を備える。EEDFなら\code{build\_pdf}、反応率なら\code{rate}、壁損失なら\code{loss\_rate}を実装する。例えば、壁損失を別の有効長モデルへ変更し、順解析と逆解析から共通して利用することができる。

現在の計算核は、外部衝突相手密度を与えた線形定常CRと、領域ごとの独立計算を採用する。この構造の中ではデータ改訂とplugin追加でモデルを展開できる。時間発展、領域間輸送、未知密度同士の積を含む化学連成は、連立方程式と解法を拡張する開発に相当する。これらの位置づけと実測検証の順序を第\ref{sec:outlook}節にまとめる。

\subsection{species packによる複数種の組立て}
species packは、状態、反応、遷移、分子帯、出典をまとめる再利用単位である。例えばArの原子データとO 777 nmの縮約データは、次の入力で一つのcaseへ組み込める。
\begin{quote}\small\ttfamily
species\_packs:\\
\hspace*{1em}- yaml\_file: data/species\_packs/ar\_2p1\_2p6\_nist.yaml\\
\hspace*{1em}- yaml\_file: data/species\_packs/o\_777\_reduced.yaml\\
\hspace*{2em}namespace: oxygen
\end{quote}
この構成では、同じ領域のArとOが共通の電子密度とEEDFを使い、それぞれの反応係数、供給密度、放射分岐から発光を計算する。Oに対するO$_2$消光のように、混合相手を反応の衝突相手として指定することもできる。

pack内の相対パスはpack自身を基準に解決する。\code{namespace}はpack内で宣言された状態・反応等のIDとその内部参照へ接頭辞を付ける。pack外の供給粒子・衝突相手は共通密度キーを参照するため、異なるpackを同じ組成条件へ接続できる。重複IDは入力時に検出する。新版への置換は旧packの選択を外して新版を指定する手順とし、使用したデータ集合を一意に定める。

\begin{table}[htbp]
\centering\small
\begin{tabularx}{\linewidth}{p{25mm}YY}
\toprule
対象例 & 現在利用できる構成 & 適用条件を整えるための入力 \\
\midrule
Ar & 2p1/2p6の放射分岐pack、六上準位の直接励起pack & 対象条件に応じた準安定密度、cascade・消光・捕獲の採否 \\
O/O$_2$ & O 777 nmの励起とO$_2$消光、Arとの合成例 & O密度と、対象条件に対応する状態別励起断面積 \\
Cl/Cl$_2$ & Cl有効発光、Arとの共有状態例、Cl$_2$帯形 & Cl供給密度、有効励起曲線、822 nmの解離励起寄与 \\
F/NF$_3$ & F線と分子帯を含む多視線の生成例 & F密度、帯の有効供給密度、励起・消光のデータ \\
CF$_4$/O$_2$/Ar & Ar/O packとCF、CF$_2$等の帯の合成例 & 各ラジカル密度と帯の係数・収率・profile \\
BCl$_3$/Ar等 & 追加caseを記述する設定ひな形 & 状態・反応・断面積・放射データの具体化 \\
\bottomrule
\end{tabularx}
\caption{現在のガス別構成と利用時の準備。Arの六上準位には外部断面積候補、O/Cl等の例には例示・有効データを用いる。検証結果は第\ref{sec:benchmarks}節、実測への展開は第\ref{sec:outlook}節に示す。}
\label{tab:gases}
\end{table}

混合ガスの計算では、親ガス分率、ラジカル・準安定密度、EEDFを入力として与える。各種の密度は一つの入力欄で定義し、解く準位密度と固定供給密度の対応を明示する。ガス添加による電子状態や解離組成の変化は、条件ごとにこの入力を更新して表現する。外部の電子・化学モデルを用いる場合も、その結果を領域ごとのEEDFと供給密度へ対応づけられる。したがって、同じ分光計算を使いながら、実測に基づく入力と外部計算に基づく入力を比較できる。

\subsection{繰返し計算と結果の再現性}
compile時には、状態参照、格子、反応構造、固有線形状、視線重み、装置の観測計画を準備する。繰返しfitでは状態量に依存する率と準位密度を更新し、固定構造を再利用する。この分離によって、物理項を振る感度解析や多数の順解析を伴う推定を同じ計算経路で実行できる。

入力の正規化、相対パス補正、schemaと意味検査、観測IDと予測の対応、校正基準の照合が、モデルと測定の接続を支える。断面積のhashと成分別の出力単位を保存し、観測、予測、残差をCSVへ出力するため、第三者は採用したデータと結果を対応づけて確認できる。通常の最適化ではスペクトル計算を反復し、最終解の詳細評価を\code{--audit}で追加する。これにより、推定に必要な反復と、結果の解釈に必要な診断を段階的に実行できる。
